\documentclass[ ../main.tex]{subfiles}
\providecommand{\mainx}{..}

\begin{document}
\section{A practical two-level perfect hash function}

\COMMENT{
\begin{algorithm}
	\caption{Two-level perfect hash function constructor}
	\label{alg:2ph}
	\SetKwProg{func}{function}{}{}
	\KwIn
	{
		$\Set{X} \subseteq \Set{U}$ is the objective set,
		$r \in \SetBuilder{q \in \RatSet}{q \in 2^{-j} \land j \in \NatSet}$ is a load factor, and
		$k \in \NatSet$ is the number of entries in the intermediate hash level.
	}
	\KwOut
	{
		A two-level perfect hash function of $\Set{X} \subseteq \Set{U}$ with a load factor $r$ and an intermediate level of $k$ indices, denoted by $\hash_{\Set{X}}^{r} \colon \Set{U} \mapsto \left\{0,\ldots,N\right\}$ where $N \coloneqq \frac{\Card{\Set{X}}}{r}$.
	}
	\func{\makehash{$\Set{X}$, $r$, $k$}}
	{
		$t \coloneqq \lceil \log_2 k \rceil$\;
		$N \coloneqq \frac{\Card{\Set{X}}}{r}$\;
		\tcp{$\Set{X}[0],\ldots,\Set{X}[k-1]$ is a total partition of $\Set{X}$.}
		$\Set{X}[\ell] \coloneqq \SetBuilder{x \in \Set{X}}{\trunc\!\left(\hash(x' \cat 0'),t\right)' \mod k = \ell}$\;
		$\Set{Y} \gets \EmptySet$\;
		\For{$\ell \in \{0,\ldots,k-1\}$}
		{
			\For{$n \gets 1$ \KwTo $\infty$}
			{
				$\Set{Y}[\ell] \gets \EmptySet$\;
				$\alpha \gets \True$\;
				\For{$x \in \Set{X}[\ell]$}
				{
					$\beta \coloneqq \trunc\!\left(\hash(x' \cat n'),N\right)$\;
					\If{$\beta \in \SetUnion[\Set{Y}[\ell]][\Set{Y}]$}
					{
						$\alpha \gets \False$\;
					}
					$\Set{Y}[\ell] \gets \SetUnion[\Set{Y}[\ell]][\{\beta\}]$\;
				}
				\If{$\alpha$}        
				{
					$\Set{Y} \gets \SetUnion[\Set{Y}][\Set{Y}[\ell]]$\;
					\Break\;
				}
			}
			%$A_\ell \gets n'$\;
			$\Set{A} \gets \SetUnion[\Set{A}][\{\Tuple{\ell,n}\}]$\;
		}
		\tcp{A space-optimal serialization of $\Tuple{\Set{A},k,N}$ is $(\ln 2)^{-1}$ bits per element if $r$ is $1$.}
		\Return $\Tuple{\Set{A},k,N}'$
	}
\end{algorithm}
}


\begin{algorithm}
	\caption{Two-level perfect hash function constructor}
	\label{alg:2ph}
	\SetKwProg{func}{function}{}{}
	\KwIn
	{
		$\Set{X} \subseteq \Set{U}$ is the objective set, $r \in \SetBuilder{q \in \RatSet}{q \in 2^{-j} \land j \in \NatSet}$ is the load factor, and $k \in \NatSet$ is the number of entries in the intermediate hash level.
	}
	\KwOut
	{
		A two-level perfect hash function of $\Set{X} \subseteq \Set{U}$ with a load factor $r$ and an intermediate level of $k$ indices, denoted by $\hash_{\Set{X}}^{r} \colon \Set{U} \mapsto \left\{0,\ldots,N\right\}$ where $N \coloneqq \frac{\Card{\Set{X}}}{r}$.
	}
	\func{\makehash{$\Set{X}$, $r$, $k$}}
	{
		$N \coloneqq \frac{\Card{\Set{X}}}{r}$\;
		$t \coloneqq \lceil \log_2 k \rceil$\;
		\tcp{$\Set{X}[1],\ldots,\Set{X}[k]$ is a total partition of $\Set{X}$ and $\Set{X}[j_1],\ldots,\Set{X}[j_k]$ is in decreasing order of cardinality.}	
		$\Set{X}[\ell] \coloneqq \SetBuilder{x \in \Set{X}}{\trunc\!\left(\hash(x' \cat 0'),t\right)' \mod k = \ell}$\;
		$\Set{Y} \gets \EmptySet$\;
		$\Set{A} \gets \EmptySet$\;		
		\For{$\ell \gets 1$ \KwTo $k$}
		{
			\For{$n \gets 1$ \KwTo $\infty$}
			{
				$\beta(x) \coloneqq \trunc\!\left(\hash(x' \cat n'),N\right)$\;
				$\Set{Y}[\ell] \coloneqq \SetBuilder{\beta(x)}{x \in \Set{X}[\ell]}$\;
				\If{$\Card{\Set{Y}[\ell]} = \Card{\Set{X}[\ell]} \land \Set{Y}[\ell] \cap \Set{Y} = \EmptySet$}
				{
					$\Set{Y} \gets \SetUnion[\Set{Y}][\Set{Y}[\ell]]$\;
					\Break\;
				}
			}
			$\Set{A} \gets \SetUnion[\Set{A}][\{\Tuple{\ell,n}\}]$\;
		}
		\Return $\Tuple{\Set{A},k,N}$
	}
\end{algorithm}





\subsection{Analysis}
Suppose we have a set $\Set{X}$ of $m$ elements with some total order $x_{(1)},\ldots,x_{(m)}$ for which we seek a perfect hash function $\hash_{\Set{X}} \colon \Set{U} \mapsto \IntSet$.

We denote the \emph{serialization} of some value $x$ by $x'$.
If $x$ is already understood to be a serialization, then $x'$ denotes deserialization instead.

By the assumption that the hash function $\hash \colon \Set{U} \mapsto \BitSet^k$ is a random oracle restricted to the codomain $\BitSet^k$, the $n$-th attempt (trial) $\hash\!\left(x_{(j)}' \cat n'\right)$ to find a non-colliding hash for $x_{(j)}$ has a probability of success $p_j = \frac{m-j+1}{m}$ since we have already found hashes for the previous $j-1$ elements and therefore there are only $m-(j-1)$ hashes remaining.

This is an example of the Coupon collector's problem.
Let $\RV{T_j}$ denote the uncertain number of trials needed to find a non-colliding hash for $x_{(j)}$.
The total number of trials needed is an uncertain random variable given by
\begin{equation}
\RV{T} = \sum_{j=1}^{m} \RV{T_j}\,.
\end{equation}
By the linearity of the expectation operator,
\begin{equation}
\Expect{\RV{T}} = \sum_{j=1}^{m} \Expect{\RV{T_j}} = \sum_{j=1}^{m} \frac{1}{p_j} = m H_{m-1}\,,
\end{equation}
which asymptotically converges to $\Expect{\RV{T}} = m \gamma + m \ln(m-1)$ or on average $\gamma + \ln (m-1)$ trials per element of $\Set{X}$.

If we use the minimum number of bits

The variance of $\RV{T}$ is given by
\begin{equation}
\Var{\RV{T}} = \sum_{j=2}^{m }\Var{\RV{T_j}}\,.
\end{equation}



\end{document}